% GENERATED FILE. Edit canonical lessons, then run npm run generate:books.
% canonical-source-hash: fnv1a64:4785f7a8e95bc3a7
% canonical-lessons: ZH-C354-tiwen, ZH-C354-huati, ZH-W354-ti-topic, ZH-C354-timu, ZH-C354-neirong, ZH-C354-tanhua

\chapter{Changing the Subject}
\label{ch:zh-changing-the-subject}
\hlchaptermodality{354}

\begin{chapteropening}
\textbf{By the end of this chapter:} \emph{I can name a topic, put a question in class and say what something is about.}

Complete the last lesson of chapter 354: I can name a topic, put a question in class and say what something is about.
\end{chapteropening}

\section[tíwèn]{tíwèn — to put a question (in class or a meeting)}

\label{lesson:ZH-C354-tiwen}



\begin{quote}
Before the new one: say the Mandarin for to argue back, then the Mandarin for to share someone's view.
\end{quote}

\subsection*{You'll want to know: tíwèn}
\emph{tíwèn} — \textquotedblleft{}to put a question (in class or a meeting)\textquotedblright{}.

\begin{quote}
\emph{yǒu wèntí zuìhǎo xiàkè yǐhòu tíwèn} — If you have a question, it's best to ask after class.
\end{quote}

Its characters come later; for now it is a word for the ear.

\begin{cousinweb}
\emph{tíwèn} is built from \emph{tí} and \zh{问} 'to ask', both inherited from Middle Chinese, \emph{dej} and \emph{mjunH} (Guangyun, in Baxter's transcription), which glosses the first 'to carry by the hand' and the second 'to inquire'.
\end{cousinweb}

\begin{grammarlens}[title={What you've built}]
One more word for giving reasons.
\end{grammarlens}

\subsection*{Your turn}
\begin{itemize}
\raggedright
  \item \emph{Say it:} \emph{tíwèn}
  \item \emph{Say it:} \emph{tíwèn}, once more
  \item \emph{Say it:} \emph{yǒu wèntí zuìhǎo xiàkè yǐhòu tíwèn} — If you have a question, it's best to ask after class.
  \item \emph{Recall:} say the Mandarin for to argue back, then the Mandarin for to share someone's view, then say \emph{tíwèn} again
\end{itemize}

\begin{tcolorbox}[breakable,colback=teal!4,colframe=teal!35!black,title={Before you move on}]
What does tíwèn mean? (\textquotedblleft{}To put a question (in class or a meeting)\textquotedblright{}.) Say it once more.
\end{tcolorbox}
\section[\texorpdfstring{\zh{话题} (huàtí)}{话题 (huàtí)}]{\zh{话题} (huàtí) — topic of conversation}

\label{lesson:ZH-C354-huati}



\begin{quote}
Before the new one: say the Mandarin for to share someone's view, then the Mandarin for to put a question.
\end{quote}

\subsection*{You'll want to know: \zh{话题}}
\textbf{\zh{话题}} — \emph{huàtí} — \textquotedblleft{}topic of conversation\textquotedblright{}.

\begin{quote}
\textbf{\zh{我觉得这个话题很有意思}} \emph{wǒ juéde zhège huàtí hěn yǒu yìsi} — I think this is an interesting topic.
\end{quote}

The new character, in pieces: \zh{题} — fifteen strokes in two pieces: on the left \zh{是} \emph{shì}, nine strokes, its last stroke sweeping out long under the right side; on the right the page piece, six strokes: a bar, a short left-falling stroke, a narrow box, then a left-falling stroke and a dot under it. The next lesson puts it in your hand.

\begin{cousinweb}
\textbf{\zh{话题}} (\emph{huàtí}) is built from \zh{话} 'talk' and \zh{题} 'a heading', both inherited from Middle Chinese, \emph{hwaejH} and \emph{dej} (Guangyun, in Baxter's transcription), which glosses the first 'speech' and the second 'the title of a book' and, quoting an older dictionary, 'the forehead'.
\end{cousinweb}

\begin{grammarlens}[title={What you've built}]
One more word for giving reasons.
\end{grammarlens}

\subsection*{Your turn}
\begin{itemize}
\raggedright
  \item \emph{Say it:} \emph{huàtí}
  \item \emph{Say it:} \emph{huàtí}, once more
  \item \emph{Say it:} \emph{wǒ juéde zhège huàtí hěn yǒu yìsi} — I think this is an interesting topic.
  \item \emph{Recall:} say the Mandarin for to share someone's view, then the Mandarin for to put a question, then say \emph{huàtí} again
\end{itemize}

\begin{tcolorbox}[breakable,colback=teal!4,colframe=teal!35!black,title={Before you move on}]
What does \zh{话题} mean? (\textquotedblleft{}Topic of conversation\textquotedblright{}.) Say it once more.
\end{tcolorbox}
\section[\texorpdfstring{\zh{题} (tí)}{题 (tí)}]{\zh{题} — shì, then the page piece}

\label{lesson:ZH-W354-ti-topic}



\begin{quote}
The word first, then the character inside it.

\begin{itemize}
\raggedright
  \item \emph{Recall:} say \emph{huàtí} — \textbf{R1}, one lesson back, before it decays
\end{itemize}

\begin{quote}
\zh{是} \zh{话} \zh{问}
\end{quote}

\emph{Read it:} each one aloud, with nothing said first — \emph{shì}, \emph{huà}, \emph{wèn}
\end{quote}

\subsection*{Script — shì, then the page piece}
\hlblockfigure{\detokenize{figures/ZH-W354-ti-topic-filmstrip.pdf}}{How \zh{题} is written, stroke by stroke}

\begin{quote}
\zh{题}
\end{quote}

\textbf{Two pieces, left before right.} Count the pieces, not the strokes.

\textbf{The left piece, \zh{是}, nine strokes}, narrower than usual: the sun box on top, a long horizontal, then the lower part, whose last stroke sweeps out long to the right, under the right piece.

\textbf{The right piece, the page piece, six strokes}: a horizontal, a short stroke falling to the left, a vertical, a horizontal that turns down, a left-falling stroke inside, then a dot at the foot.

9 + 6 = fifteen, and the first nine your hand already knows.

\subsection*{Your turn}
\begin{itemize}
\raggedright
  \item \emph{Write it:} \zh{题}, all fifteen, with the model in front of you
  \item \emph{Cover:} write it again from memory, and count the pieces, not the strokes
\end{itemize}

\begin{tcolorbox}[breakable,colback=teal!4,colframe=teal!35!black,title={Before you move on}]
How many pieces in \zh{题}? (\textbf{Two}.) What is the left piece? (\textbf{\zh{是}}.) What does its last stroke do? (\textbf{It sweeps out long, under the page piece}.)

Source: \href{https://raw.githubusercontent.com/chanind/hanzi-writer-data/68d10a4b21150cae5e1ebbd223eed289cf32d90c/data/\%E9\%A2\%98.json}{Hanzi Writer Data \zh{题}, pinned stroke paths and medians}.
\end{tcolorbox}
\section[tímù]{tímù — title; question on a test}

\label{lesson:ZH-C354-timu}



\begin{quote}
Before the new one: say the Mandarin for to put a question, then the Mandarin for a topic of conversation.
\end{quote}

\subsection*{You'll want to know: tímù}
\emph{tímù} — \textquotedblleft{}title; question on a test\textquotedblright{}.

\begin{quote}
\emph{wǒ kàn bu dǒng zhège tímù} — I can't make sense of this question.
\end{quote}

Its characters come later; for now it is a word for the ear.

\begin{cousinweb}
\emph{tímù} is built from \zh{题} 'a heading' and \emph{mù} 'the eye', both inherited from Middle Chinese, \emph{dej} and \emph{mjuwk} (Guangyun, in Baxter's transcription), which glosses the first 'the title of a book' and, quoting an older dictionary, the second 'the human eye'.
\end{cousinweb}

\begin{grammarlens}[title={What you've built}]
One more word for giving reasons.
\end{grammarlens}

\subsection*{Your turn}
\begin{itemize}
\raggedright
  \item \emph{Say it:} \emph{tímù}
  \item \emph{Say it:} \emph{tímù}, once more
  \item \emph{Say it:} \emph{wǒ kàn bu dǒng zhège tímù} — I can't make sense of this question.
  \item \emph{Recall:} say the Mandarin for to put a question, then the Mandarin for a topic of conversation, then say \emph{tímù} again
\end{itemize}

\begin{tcolorbox}[breakable,colback=teal!4,colframe=teal!35!black,title={Before you move on}]
What does tímù mean? (\textquotedblleft{}Title; question on a test\textquotedblright{}.) Say it once more.
\end{tcolorbox}
\section[nèiróng]{nèiróng — content; what something is about}

\label{lesson:ZH-C354-neirong}



\begin{quote}
Before the new one: say the Mandarin for a topic of conversation, then the Mandarin for a question on a test.
\end{quote}

\subsection*{You'll want to know: nèiróng}
\emph{nèiróng} — \textquotedblleft{}content; what something is about\textquotedblright{}.

\begin{quote}
\emph{zhè běn shū de nèiróng hěn yǒu yìsi} — This book has interesting content.
\end{quote}

Its characters come later; for now it is a word for the ear.

\begin{cousinweb}
\emph{nèiróng} is built from \emph{nèi} and \emph{róng}, both inherited from Middle Chinese, \emph{nwojH} and \emph{yowng} (Guangyun, in Baxter's transcription), which glosses the first 'to enter' and the second 'to hold'.
\end{cousinweb}

\begin{grammarlens}[title={What you've built}]
One more word for giving reasons.
\end{grammarlens}

\subsection*{Your turn}
\begin{itemize}
\raggedright
  \item \emph{Say it:} \emph{nèiróng}
  \item \emph{Say it:} \emph{nèiróng}, once more
  \item \emph{Say it:} \emph{zhè běn shū de nèiróng hěn yǒu yìsi} — This book has interesting content.
  \item \emph{Recall:} say the Mandarin for a topic of conversation, then the Mandarin for a question on a test, then say \emph{nèiróng} again
\end{itemize}

\begin{tcolorbox}[breakable,colback=teal!4,colframe=teal!35!black,title={Before you move on}]
What does nèiróng mean? (\textquotedblleft{}Content; what something is about\textquotedblright{}.) Say it once more.
\end{tcolorbox}
\section[tánhuà]{tánhuà — to have a serious talk}

\label{lesson:ZH-C354-tanhua}



\begin{quote}
Before the new one: say the Mandarin for a question on a test, then the Mandarin for what something is about.
\end{quote}

\subsection*{You'll want to know: tánhuà}
\emph{tánhuà} — \textquotedblleft{}to have a serious talk\textquotedblright{}.

\begin{quote}
\emph{lǎoshī xiǎng gēn nǐ tánhuà} — The teacher wants a word with you.
\end{quote}

Its characters come later; for now it is a word for the ear.

\begin{cousinweb}
\emph{tánhuà} is built from \emph{tán} and \zh{话} 'talk', both inherited from Middle Chinese, \emph{dam} and \emph{hwaejH} (Guangyun, in Baxter's transcription), which glosses the first 'to talk' and the second 'speech'.
\end{cousinweb}

\begin{grammarlens}[title={What you've built}]
One more word for giving reasons.
\end{grammarlens}

\subsection*{Your turn}
\begin{itemize}
\raggedright
  \item \emph{Say it:} \emph{tánhuà}
  \item \emph{Say it:} \emph{tánhuà}, once more
  \item \emph{Say it:} \emph{lǎoshī xiǎng gēn nǐ tánhuà} — The teacher wants a word with you.
  \item \emph{Recall:} say the Mandarin for a question on a test, then the Mandarin for what something is about, then say \emph{tánhuà} again
\end{itemize}

\begin{tcolorbox}[breakable,colback=teal!4,colframe=teal!35!black,title={Before you move on}]
What does tánhuà mean? (\textquotedblleft{}To have a serious talk\textquotedblright{}.) Say it once more.
\end{tcolorbox}
